\section{Overlapping comparison windows}\label{sec:windows}

The comparison \(\Comp^0\) dominates the syndrome indefinitely by
\eqref{eq:domination}, but its full backward cone eventually wraps.
We now use only short, clean-start comparison estimates to control
that same old field.

Fix all parameters from Proposition~\ref{prop:path}. For sufficiently
large \(M\), set
\begin{equation}\label{eq:epoch}
 h=\left\lfloor\frac{M}{100(R+3)}\right\rfloor,\qquad
 a_j=jh\quad(j\ge0).
\end{equation}
An epoch consists of \(h\) slabs, and a two-epoch comparison window
has depth \(2h\). It satisfies \(2R(2h)<M/2\), so all results of
Section~\ref{sec:contours} apply within it. We assume \(h\ge1\).

On one common sample of physical marks let
\(\Comp_j=\Comp^{a_j}\). Each of these auxiliary fields starts all
ones at its own date; none restarts the device. Define
\begin{align}
 G_j&=\{\Comp_j(z,a_{j+2})=\Comp_{j+1}(z,a_{j+2})
                       \text{ for every }z\},\label{eq:G}\\
 V_j&=\left\{\begin{array}{l}
 \text{at no date }k\in(a_j,a_{j+2}]
 \text{ does }\Proxy^{a_j}(\cdot,k)\\
 \text{contain a spatially connected set of diameter }\ge M/200
 \end{array}\right\}.\label{eq:V}
\end{align}
Spatial adjacency here is sup distance one. The diameter is the
ambient coarse torus diameter.

\begin{lemma}[Coalescence over an overlap]\label{lem:coalescence}
There are constants \(C,\kappa>0\), uniform in \(j,L\) and the
allowed rates, such that
\begin{equation}\label{eq:window-cost}
 \Prb(G_j^c)\le CM^4e^{-\kappa h},\qquad
 \Prb(V_j^c)\le C(2h)M^4e^{-\kappa M/200}.
\end{equation}
Both estimates hold conditionally on the entire physical history
through \(a_jT\).
\end{lemma}
\begin{proof}
During the common part of their lifetimes,
\(\Comp_j\le\Comp_{j+1}\): the latter is initialized all ones at
\(a_{j+1}\), and every later map is monotone or constant.
Suppose the two fields disagree at a specified endpoint
\((x,a_{j+2})\). The older bit there is zero and the newer bit is
one. They use the same update map. A constant map cannot give
different outputs. For a nonconstant map, different outputs imply
that some predecessor argument differs, with displacement in
\(\Delta\). Choose one such predecessor and repeat backwards.
This produces a path of disagreements at every integer date back
to \(a_{j+1}\), including that date.

Every site on this path is a zero of \(\Comp_j\). Its dates lie
strictly after \(a_j\), since \(a_{j+1}>a_j\), so every site belongs
to \(\Proxy^{a_j}\). Consecutive sites have product distance at most
\(R\). The endpoint-to-endpoint temporal displacement is \(h\).
Proposition~\ref{prop:path}, anchored at the specified final site,
bounds this event by \(Ce^{-\kappa h}\) in the two-epoch window.
Summing over the \(M^4\) possible final spatial sites proves the
bound for \(G_j^c\).

If \(V_j\) fails, at some one of its \(2h\) dates a spatial path in
the proxy joins two sites separated by at least \(M/200\).
The finite coarse torus ensures that a connected set witnessing
the diameter contains such a finite path. Apply
Proposition~\ref{prop:path} at its initial spatial site, with
\(u=M/200\), and sum over \(M^4\) sites and \(2h\) dates. This proves
the other estimate.

Both comparisons in this argument are clean-start functions of the
increments after \(a_jT\). Those increments are independent of the
history through \(a_jT\). Thus precisely the same estimates hold
under the stated conditioning. No conditional estimate is used at
the middle date of the overlap.
\end{proof}

\begin{lemma}[Propagation of comparison agreement]\label{lem:iteration}
For every integer \(n\ge1\), on
\[
 \bigcap_{0\le j<n}(G_j\cap V_j)
\]
the physical syndrome, started empty, has only components of physical
diameter less than \(L/100\) at every time in \([0,nhT]\).
\end{lemma}
\begin{proof}
We track which clean-start comparison equals the old comparison.
At the beginning of epoch \(i\ge1\) we claim
\begin{equation}\label{eq:seam}
 \Comp^0(\cdot,a_i)=\Comp_{i-1}(\cdot,a_i).
\end{equation}
For \(i=1\) this is identity, since \(\Comp_{i-1}=\Comp_0\).
If it holds at \(a_i\), the two fields use common deterministic
maps and remain equal through the following epoch. At \(a_{i+1}\),
the event \(G_{i-1}\) equates \(\Comp_{i-1}\) with \(\Comp_i\).
This proves \eqref{eq:seam} inductively for all needed \(i\).

Let \(W_k\) be the blocks visited by syndrome vertices at any instant
of slab \(I_k\). Equation~\eqref{eq:domination} puts \(W_k\) in
\(\Proxy^0(\cdot,k)\). In epoch zero this is directly controlled by
\(V_0\). In epoch \(i\ge1\), equality
\eqref{eq:seam} and its propagation ensure equality of the bits
used to form \(\Proxy^0(\cdot,k)\) and
\(\Proxy^{a_{i-1}}(\cdot,k)\). This includes the predecessor date of
the earliest slab in the epoch: it is exactly \(a_i\).
Thus \(W_k\) is contained in the proxy tested by \(V_{i-1}\).
The physical syndrome at that seam can be nonempty; only comparison
bits have been equated.

Two adjacent physical vertices belong to identical or adjacent
blocks. Further, if their containing blocks have coarse torus
distance \(d\), their physical torus distance is at most \(b(d+1)\).
To verify this inequality, lift their relative block displacement
in each coordinate by a shortest representative; their offsets
inside the two intervals add less than one block width. This works
also for noninteger \(b=L/M\).

If a physical component had diameter at least \(L/100\), its
edge-connected block image would therefore have diameter at least
\[
 M/100-1\ge M/200
\]
for sufficiently large \(M\). That image is contained in the
visited set in its slab, contradicting the applicable \(V_j\).
The whole-slab definition covers both configurations immediately
before and after every physical jump. Almost surely no ring
occurs at a deterministic slab boundary. At time zero the syndrome
is empty. These observations give the assertion for all times in
the stated interval.
\end{proof}

\begin{proposition}[Cumulative finite-torus estimate]\label{prop:history}
There exist \(\eta_0,C_0,\beta>0\) and an integer \(L_0\), independent
of the allowed spatial noise rates, such that in either sector,
for every \(L\ge L_0\) and deterministic \(t\ge0\),
\begin{equation}\label{eq:history}
 \Prb\{\exists s\in[0,t]:
       \text{a syndrome component at }s
                   \text{ has diameter }\ge L/100\}
 \le \min\{1,C_0(1+t)e^{-\beta L}\}.
\end{equation}
The initial syndrome is empty.
\end{proposition}
\begin{proof}
For \(t>0\) put \(n=\lceil t/(hT)\rceil\). A union bound over the
events in Lemma~\ref{lem:iteration}, with
Lemma~\ref{lem:coalescence}, gives
\begin{equation}\label{eq:polynomial-history}
 CnM^4\bigl[e^{-\kappa h}
                  +2h e^{-\kappa M/200}\bigr]
\end{equation}
as an upper bound on the probability. Some of these tests extend
past \(t\). Using their future marks in an upper bound is legitimate;
neither the device nor the recovery consults them.

For all sufficiently large \(M\),
\(h\ge M/[200(R+3)]\), while \(h\le M\).
Also \(M\ge L/(2B)\) once \(L\) is sufficiently large.
Consequently the bracket in
\eqref{eq:polynomial-history} is bounded by a polynomial in \(L\)
times \(e^{-b_1L}\) for a fixed \(b_1>0\). Since
\(hT\ge1\), we have \(n\le1+t\). Absorb all remaining polynomial
factors into half the positive exponential rate, enlarging \(L_0\)
and \(C_0\). This gives \eqref{eq:history} with, for example, some
\(\beta\le b_1/2\). The event at \(t=0\) is empty.
The probability is always at most one.

Overlapping windows share many increments. Their independence was
not used: the entire accumulation step is a union bound on the
original sample of marks. The argument also never assumes
contraction for arbitrary comparison data at a restart.
\end{proof}
